% Generated draft: requires mathematical and source review.
\documentclass[11pt]{article}
\usepackage[margin=25mm]{geometry}
\usepackage{fontspec}
\setmainfont{Latin Modern Roman}
\usepackage{amsmath,amssymb,mathtools,unicode-math}
\setmathfont{Latin Modern Math}
\usepackage{enumitem,needspace,xurl,hyperref}
\setlength{\parindent}{0pt}
\setlength{\parskip}{5pt}
\setlength{\emergencystretch}{3em}
\newcommand{\N}{\mathbb{N}}
\newcommand{\Z}{\mathbb{Z}}
\newcommand{\Q}{\mathbb{Q}}
\newcommand{\R}{\mathbb{R}}
\newcommand{\C}{\mathbb{C}}
\newcommand{\F}{\mathbb{F}}
\newcommand{\A}{\mathbb{A}}
\newcommand{\PP}{\mathbb{P}}
\newcommand{\E}{\mathbb{E}}
\newcommand{\eps}{\varepsilon}
\newcommand{\dd}{\,\mathrm{d}}
\begin{document}
\section*{UnsolvedMath ID: 800092}
\textbf{Source problem code:} catalogue entry
\subsection{Definitions and mathematical statement}
Let \(p\colon M\to B\) be a fibration with fiber \(F\), and let \(\mathcal F\) be the foliation defined by this fibration. The base \(B\) is closed, meaning compact and without boundary under the usual manifold convention, and the source assumes
\[
\pi_1(F)\cong\mathbb Z^k.
\]
Here \(\mathbb Z^k\) denotes the free abelian group of rank \(k\). For \(b\in B\), write \(F_b=p^{-1}(b)\).

The question is to identify additional hypotheses on this setting under which every foliation \(\mathcal G\) sufficiently close to \(\mathcal F\) has a compact leaf close to a fiber. Schematically, the desired conclusion is
\[
\text{for every sufficiently close }\mathcal G,\quad
\exists b\in B\ \exists L\subseteq M:\quad
L\text{ is a compact leaf of }\mathcal G
\text{ and is close to }F_b.
\]
This is a request for hypotheses yielding such a persistence statement, rather than an assertion that the displayed assumptions alone suffice.

\textbf{Unresolved qualifications.} The source does not specify the regularity of the fibration or foliations, the topology measuring closeness of foliations, or the meaning and quantitative strength of closeness of a leaf to a fiber. It also does not state whether a fiber is prescribed in advance or may depend on \(\mathcal G\); the schematic display uses the latter reading only to express the existential wording. Compactness, connectedness, and boundary conditions for \(F\), and the intended range of \(k\), are not explicitly supplied. If fibers need not be connected, the convention relating fibers to leaves must also be specified. These choices must be resolved before the question becomes a fully quantified mathematical statement.
\subsection{Short English statement}
Given a foliation arising from a fibration over a closed base, with fiber fundamental group \(\mathbb Z^k\), determine additional hypotheses ensuring that every sufficiently close foliation has a compact leaf close to a fiber. The source leaves both notions of closeness and several geometric assumptions unspecified.
\subsection{Sources}
\begin{itemize}
\item Local UnsolvedMath registry: permanent ID \(800092\), exact source problem code \texttt{catalogue entry}, title ``Foliations and Diffeomorphism Groups --- Question 33.'' The record attributes the problem to Nariman and Yazdi, with proposed year 2025.
\item Nariman, Yazdi, \emph{Problem List on Foliations and Diffeomorphism Groups} (2025), Question 33. Supplied source URL: \url{https://arxiv.org/abs/2503.18274}. The supplied provenance records access on 2026-10-16 and extraction by \texttt{arxiv\_latex}; neither access nor extraction has been independently verified here.
\item The supplied registry labels this problem ``open'' and records an update on 2026-10-16. This is a supplied status label, not an independently verified mathematical or literature assessment.
\end{itemize}
\end{document}
